\begin{lemma}
\label{lemma-ML-countable}
Let $R$ be a ring.
Let $M$ be an $R$-module.
Assume $M$ is Mittag-Leffler and countably generated.
For any $R$-module map $f : P \to M$ with $P$ finitely generated there
exists an endomorphism $\alpha : M \to M$ such that
\begin{enumerate}
\item $\alpha : M \to M$ factors through a finitely presented $R$-module, and
\item $\alpha \circ f = f$.
\end{enumerate}
\end{lemma}

\begin{proof}
Write $M = \colim_{i \in I} M_i$ as a directed colimit of finitely
presented $R$-modules with $I$ countable, see
Lemma \ref{lemma-ML-countable-colimit}.
The transition maps are denoted $f_{ij}$ and we use $f_i : M_i \to M$
to denote the canonical maps into $M$. Set $N = \prod_{s \in I} M_s$. Denote
$$
M_i^* = \Hom_R(M_i, N) = \prod\nolimits_{s \in I} \Hom_R(M_i, M_s)
$$
so that $(M_i^*)$ is an inverse system of $R$-modules over $I$.
Note that $\Hom_R(M, N) = \lim M_i^*$.
As $M$ is Mittag-Leffler, we find for every
$i \in I$ an index $k(i) \geq i$ such that
$$
E_i := \bigcap\nolimits_{i' \geq i} \Im(M_{i'}^* \to M_i^*)
=
\Im(M_{k(i)}^* \to M_i^*)
$$
Choose and fix $j \in I$ such that $\Im(P \to M) \subset \Im(M_j \to M)$.
This is possible as $P$ is finitely generated. Set $k = k(j)$.
Let
$x = (0, \ldots, 0, \text{id}_{M_k}, 0, \ldots, 0) \in M_k^*$ and
note that this maps to $y = (0, \ldots, 0, f_{jk}, 0, \ldots, 0) \in M_j^*$.
By our choice of $k$ we see that $y \in E_j$. By
Example \ref{example-ML-surjective-maps}
the transition maps $E_i \to E_j$ are surjective for each $i \geq j$
and $\lim E_i = \lim M_i^* = \Hom_R(M, N)$. Hence
Lemma \ref{lemma-ML-limit-nonempty}
guarantees there exists an element $z \in \Hom_R(M, N)$
which maps to $y$ in $E_j \subset M_j^*$. Let $z_k$ be the $k$th component
of $z$. Then $z_k : M \to M_k$ is a homomorphism such that
$$
\xymatrix{
M \ar[r]_{z_k} & M_k \\
M_j \ar[ru]_{f_{jk}} \ar[u]^{f_j}
}
$$
commutes. Let $\alpha : M \to M$ be the composition
$f_k \circ z_k : M \to M_k \to M$.
Then $\alpha$ factors through a finitely presented module by construction and
$\alpha \circ f_j = f_j$. Since the image of $f$ is contained in the image of
$f_j$ this also implies that $\alpha \circ f = f$.
\end{proof}
